\begin{figure}[!t]
\centering
\hspace*{-0.2cm}
\includegraphics[width=1.02\columnwidth]{figures/fate_overview9d.pdf}
\caption{
\label{fig:overview}
\textbf{An overview of \approach{},} which improves attribution for a text passage via \emph{Research \& Revision}. 
Given the input text passage, the \textbf{\textcolor{BrickRed}{research stage}} uses a \textit{query generator} to raise questions about different aspects of the text.
The \textit{retriever} then searches for evidence to investigate each query. The \textbf{\textcolor{DarkOrchid}{revision stage}} first runs an \textit{agreement model} to detect disagreement between the text and the evidence, then runs an \textit{edit model} to revise the text if needed.
Finally, $M$ evidence snippets are selected to form an attribution report.
}
\end{figure}

\section{Task formulation}\label{sec:task}

We propose the task of \textit{\taskName} as follows. As Figure~\ref{fig:api} shows, the input to the system is a text passage $x$ produced by a generation model. %
The output is a revised text passage $y$ along with an \emph{attribution report} $A$, which contains \emph{evidence snippets} $e_1, \dots, e_M$ that support the content in $y$.
Optionally, the attribution report can contain additional information such as the alignment between evidence snippets and relevant parts in $y$.

We propose to measure the quality of the revised text $y$ and attribution report $A$ along two dimensions: (1) \textbf{attribution}: how much of the revised text $y$ can be attributed to the evidence in $A$,
and (2) \textbf{preservation}: how much the revised text $y$ preserves aspects of the original text $x$.


\subsection{Measuring attribution}

Previously, \citet{Rashkin2021AIS} proposed \emph{Attributable to Identified Sources} (AIS), a human evaluation framework which considers a binary notion of attribution.
Roughly speaking, a text passage $y$ is attributable to a set $A$ of evidence if a generic hearer would affirm the statement ``According to $A$, $y$'' under the context of $y$.
A system either receives full credit (1.0) if \emph{all} content in $y$ can be attributed to $A$, and no credit (0.0) otherwise.

We propose a more fine-grained, sentence-level extension of AIS.
We ask annotators to give an AIS score for each sentence $s$ of $y$, and then report the average AIS score across all sentences:
\begin{equation}
\attrais(y, A) = \operatorname*{avg}_{s \in y} \text{AIS}(s, A).
\end{equation}
Since the AIS score is binary, this effectively measures the percentage of sentences in $y$ that are fully attributed to $A$. When judging each sentence, we also give annotators access to the surrounding sentences and other necessary context, such as the question that the text passage responded to.
We also impose the maximum number of evidence snippets in the attribution report $A$
to make it concise enough for both the annotator and downstream users.
By manually inspecting 30 examples from our benchmarks,
we found $M = 5$ snippets to be sufficient for full attribution.

During model development, we define
an automated metric, auto-AIS ($\attrauto{}$), that approximates human AIS judgments.
We utilize the natural language inference (NLI) model from \citet{Honovich2022TRUERF}, which correlates well with AIS scores.
For each sentence $s$ of $y$,
and for each evidence snippet $e$ in $A$,
let $\text{NLI}(e, s)$ be the model probability of $e$ entailing $s$. We then define
\begin{equation}
\attrauto(y, A) = \operatorname*{avg}_{s \in y} \, \max_{e \in A} \text{NLI}(e, s).
\end{equation}
To improve accuracy, we decontextualize~\cite{choi2021making} each sentence based on the entire context of $y$ %
before computing the scores.
See Appendix~\ref{sec:app-auto-eval} for implementation details.

\subsection{Measuring preservation}\label{sec:preservation}

To measure preservation, we first ask annotators to decide if the revision preserves the text's original intent (completely, somewhat, or not at all --- see Appendix~\ref{sec:app-human-eval} for exact rubrics).
Like AIS evaluation, we give annotators the necessary surrounding context.
We define the binary metric $\presintent(x, y)$ to be 1.0 if the revision completely preserves the original intent, and 0.0 otherwise.

However, even if a revision preserves intent, it may still make superfluous modifications, such as reordering words, changing textual style, or including unnecessary additional information \cite{thorne2021evidence}. Different tasks have different requirements for what should be preserved. Here, we desire a simple metric that can be readily computed for many tasks and that generally penalizes unnecessary changes.
We thus define a metric based on the character-level Levenshtein edit distance \cite{levenshtein1966binary} %
between $x$ and $y$:
\begin{equation}
\preslev(x, y) = \max\left(1 - \frac{\text{{Lev}}(x, y)}{\text{{length}}(x)}, 0\right)
\end{equation}
This metric is 1.0 if $x$ and $y$ are the same, and 0.0 if $y$ completely overwrites all parts of $x$. $\preslev$ is generally sensitive to any kind of change, but certainly does not capture all notions of preservation (e.g., preserving rhyme schemes or puns).

We want the revision to preserve the original intent while avoiding superfluous edits. To reflect this, we finally combine the two metrics as
\begin{equation}
\prescomb(x, y) = \presintent(x, y) \cdot \preslev(x, y).
\end{equation}
which is 0.0 if the revision changes the intent and equal to $\preslev(x, y)$ otherwise.
Since $\presintent$ requires human annotation,
we use $\preslev$ as an automated metric for model development.






\subsection{Discussion}


Optimizing for attribution alone cannot ensure a good revision: for example, an adversarial editor could ensure 100\% attribution by simply replacing the input $x$ with the text of any arbitrary retrieved document, which is trivially attributable to itself.
Ideally, we want to maximize both attribution and preservation, while navigating any tradoffs between the two.
In our experiments, we report both metrics, as well as their harmonic mean ($\HM{}$, analogous to how recall and precision are combined in F1).


We emphasize that this evaluation scheme does not require any ``gold'' or ``reference'' edits (unlike many prior evaluations of text revision models), which are often only available for specialized domains. This enables us to broaden the scope to a much wider range of generation tasks.